\section{The Fisher model with phenotypic assortment}
\label{sec:fisher-model}

A close fit to correlations between relatives is informative only to the
extent that we understand what generates the fitted formulas. The same
pattern of resemblance can support different explanations when different
transmission mechanisms are allowed. Spelling out the model therefore
serves two purposes: it identifies which conclusions follow from the
assumptions, and it makes clear which explanations have been excluded
before the data are examined. In particular, the absence of environmental
transmission must not be mistaken for an empirical finding that
environmental transmission is unimportant. We derive the additive,
dominance-free version of \citet{Fisher1918}'s model, following the organization of
\citet{Goldberger1978} and \citet{OttoEtAl1994}: individual components,
mating, and transmission. Measurement error is omitted throughout this
section.

\subsection{Individuals, parameters, and environmental restrictions}
\label{sec:model-components}

For person $i$, define the two column vectors
\[
 X_i=\underbrace{\begin{pmatrix}A_i\\E_i\end{pmatrix}}_{\text{random components}},
 \qquad
 q=\underbrace{\begin{pmatrix}h\\e\end{pmatrix}}_{\text{common parameters}},
 \qquad P_i=q^{\mathsf T}X_i=hA_i+eE_i.
\]
$A_i$ is the person's standardized additive genetic value; $E_i$ is the
standardized environmental component. $X_i$ varies across people, whereas
$q$ contains coefficients common to the population. The superscript
$\mathsf T$ turns a column into a row. These vectors are bookkeeping
devices, not extra components of the model. There are no dominance or other nonadditive genetic components, and no gene--environment interaction term. Assume
\begin{equation}
 \begin{gathered}
 \mathbb E[A_i]=\mathbb E[E_i]=0,\qquad
 \operatorname{Var}(A_i)=\operatorname{Var}(E_i)=1,\\
 \operatorname{Cov}(A_i,E_i)=0,\qquad
 0<h<1,\qquad e=\sqrt{1-h^2}.
 \end{gathered}
 \label{eq:model-raw}
\end{equation}
Thus $\operatorname{Cov}(X_i,X_i)=I_2$, the two-dimensional identity
matrix, and
\[
 \operatorname{Var}(P_i)=h^2+e^2+2he\operatorname{Cov}(A_i,E_i)=1.
\]
The contribution of genes on the phenotype scale is $hA_i$, with variance
$h^2$. Consequently $h^2$ is narrow-sense heritability in this population;
it is not $\operatorname{Var}(A_i)$, which we normalized to one.
$E_i$ is a substantive environmental component, not measurement error.

The sexes have the same marginal distributions, and sampling people as
parents does not change those distributions. We use the same parameters
and component variances across generations. These assumptions exclude
changes in the composition of the sampled population; stationarity of
the variances is checked below.

The model also excludes a common environmental component among biological
relatives. For the pedigrees derived below, distinct biological relatives
have $\operatorname{Cov}(E_i,E_j)=0$; an offspring's environment has zero
covariance with the components of its parents and with the genes and
environments of its siblings. We implement these exclusions formally in
the birth process below. They are additional family assumptions, not
consequences of phenotypic assortment.

\citet[pp.~10--11]{Goldberger1978} emphasizes the substantive severity of
this restriction: the classical model allows the premarital environments
of spouses to resemble one another through assortment, while excluding
environmental resemblance between siblings or between parents and
children. Shared rearing supplies no additional source of similarity;
being raised together rather than apart has no effect on the predicted
correlation for otherwise comparable biological relatives. Genes are the
only component transmitted in the model. This is an assumption about the
sources of familial resemblance, not evidence establishing their relative
importance. It does \emph{not}, however, set
$\operatorname{Cov}(E_{\rm parent},A_{\rm child})$ to zero: assortment
generates that covariance, as we show below.

\subsection{What phenotypic homogamy assumes}
\label{sec:model-ph}

Let $M$ denote the wife/mother and $F$ the husband/father. Define
\[
 \rho=\operatorname{Corr}(P_M,P_F),\qquad
 \rho_A=\operatorname{Corr}(A_M,A_F),\qquad 0<\rho<1.
\]
These definitions alone say nothing about the matching mechanism.
Our formal phenotype-only matching assumption is
\begin{equation}
 X_M\perp X_F\mid P_F,\qquad X_F\perp X_M\mid P_M.
 \tag{PH}\label{eq:model-ph}
\end{equation}
Once the wife's phenotype is known, learning her particular genetic and
environmental components gives no further information about the
distribution of the husband's components; the analogous statement holds
with husband and wife exchanged. The two statements are separate
restrictions. Symmetry of conditional independence exchanges the two
vectors but does not change the conditioning phenotype.

For clarity, $Z\perp W\mid V$ means that, conditional on $V$, the joint
probability of any event concerning $Z$ and any event concerning $W$
is the product of their conditional probabilities. Equivalently, learning
$W$ does not change the conditional distribution of $Z$ once $V$ is known.
All conditional statements hold wherever the conditional laws are
defined, up to events of probability zero. Independence is preserved by
taking functions: an event about a function of $Z$ is itself an event
about $Z$, so the same factorization applies.

We additionally assume the \emph{own-person} conditional mean is linear:
\begin{equation}
 \mathbb E[A_i\mid P_i]=hP_i.
 \tag{L}\label{eq:model-linearity}
\end{equation}
Here the prediction is of a person's genetic value from that same person's
phenotype, not from the spouse's phenotype. To see why the coefficient is
$h$, write a linear conditional mean as $\alpha+\beta P_i$. Zero means give
$\alpha=0$, and
\[
 \beta=\frac{\operatorname{Cov}(A_i,P_i)}{\operatorname{Var}(P_i)}
 =\frac{h\operatorname{Var}(A_i)+e\operatorname{Cov}(A_i,E_i)}{1}=h.
\]
A population linear projection always has this coefficient, but its
conditional mean need not equal that projection. Thus (L) is a real
assumption. From $eE_i=P_i-hA_i$ it follows that
\[
 \mathbb E[E_i\mid P_i]
 =\frac{P_i-h\mathbb E[A_i\mid P_i]}{e}
 =\frac{1-h^2}{e}P_i=eP_i,\qquad
 \mathbb E[X_i\mid P_i]=qP_i.
\]
Otto et al.'s Gaussian specification supplies such linearity. Joint
normality is sufficient here, but the spouse derivation needs only (L)
together with (PH) and the stated normalization.

\subsection{Deriving the restriction on spousal genetic resemblance}
\label{sec:model-spouses}

We spell out the conditioning steps because they separate a probability
identity from the assumptions that give the model its content.

For this argument, the four components $(A_M,E_M,A_F,E_F)$ are assumed
to have an ordinary joint density.
Because $e>0$, replacing each $E_i$ by $P_i=hA_i+eE_i$ is an invertible
change of variables, so $(A_M,A_F,P_M,P_F)$ also has an ordinary density.
We do not assume an ordinary joint density for $(A_i,E_i,P_i)$, which
contains a redundant variable.

Subscripts distinguish different density functions; lower-case letters
denote values of the corresponding random variables. Let
\[
 J=f_{A_M,A_F,P_M,P_F}(a_m,a_f,p_m,p_f),\quad
 D=f_{A_F,P_M,P_F}(a_f,p_m,p_f),\quad
 H=f_{P_M,P_F}(p_m,p_f).
\]
By the definition of a conditional density, at points with positive
denominators,
\begin{align}
 f_{A_M,A_F\mid P_M,P_F}
 &=\frac{J}{H}
   =\frac{J}{D}\frac{D}{H}\notag\\
 &=f_{A_M\mid A_F,P_M,P_F}\,
   f_{A_F\mid P_M,P_F}.
 \label{eq:model-multiplication}
\end{align}
Arguments on the last line are respectively
$(a_m\mid a_f,p_m,p_f)$ and $(a_f\mid p_m,p_f)$.
This multiplication rule is an identity; no independence has been used.

Now invoke (PH). Its second restriction, by symmetry and preservation
under functions, gives $A_M\perp(A_F,P_F)\mid P_M$, because
$A_M$ is an element of $X_M$ and $(A_F,P_F)$ is a function of $X_F$.
Conditional independence therefore factors the numerator in the
following definition:
\begin{align}
 f_{A_M\mid A_F,P_M,P_F}
 &=\frac{f_{A_M,A_F,P_F\mid P_M}}{f_{A_F,P_F\mid P_M}}\notag\\
 &=\frac{f_{A_M\mid P_M}\,f_{A_F,P_F\mid P_M}}
         {f_{A_F,P_F\mid P_M}}
 =f_{A_M\mid P_M}.
 \label{eq:model-wife-reduction}
\end{align}
The cancellation is valid on the conditional support. Similarly, the
first PH restriction gives $A_F\perp P_M\mid P_F$, and hence
\begin{align}
 f_{A_F\mid P_M,P_F}
 &=\frac{f_{A_F,P_M\mid P_F}}{f_{P_M\mid P_F}}\notag\\
 &=\frac{f_{A_F\mid P_F}\,f_{P_M\mid P_F}}{f_{P_M\mid P_F}}
 =f_{A_F\mid P_F}.
 \label{eq:model-husband-reduction}
\end{align}
Substitute \eqref{eq:model-wife-reduction} into the first factor of
\eqref{eq:model-multiplication}, and
\eqref{eq:model-husband-reduction} into the second:
\begin{equation}
 f_{A_M,A_F\mid P_M,P_F}(a_m,a_f\mid p_m,p_f)
 =f_{A_M\mid P_M}(a_m\mid p_m)
  f_{A_F\mid P_F}(a_f\mid p_f).
 \label{eq:model-density-factor}
\end{equation}
Thus, given both phenotypes, the genetic values are independent, and
each retains its own-phenotype conditional distribution. Independence
given both phenotypes \emph{alone} would not establish the latter fact.

The conditional product mean is consequently
\begin{align*}
 \mathbb E[A_MA_F\mid p_m,p_f]
 &=\int\!\!\int a_ma_f
 f_{A_M\mid P_M}(a_m\mid p_m)
 f_{A_F\mid P_F}(a_f\mid p_f)\,da_m\,da_f\\
 &=\left\{\int a_m f_{A_M\mid P_M}(a_m\mid p_m)\,da_m\right\}
   \left\{\int a_f f_{A_F\mid P_F}(a_f\mid p_f)\,da_f\right\}\\
 &=\mathbb E[A_M\mid p_m]\mathbb E[A_F\mid p_f]
   =(hp_m)(hp_f).
\end{align*}
The factorization uses PH; the final equality uses (L).
Finite second moments justify these product expectations. Finally,
average over the \emph{joint spousal phenotype distribution}, using the
law of iterated expectations:
\begin{align}
 \rho_A
 &=\operatorname{Cov}(A_M,A_F)
   =\mathbb E[A_MA_F]\notag\\
 &=\mathbb E_{(P_M,P_F)}
       \{\mathbb E[A_MA_F\mid P_M,P_F]\}\notag\\
 &=h^2\mathbb E_{(P_M,P_F)}[P_MP_F]
   =\boxed{h^2\rho}.
 \label{eq:m-def}
\end{align}
The first line uses zero means and unit genetic variances; the last uses
zero means and unit phenotype variances. The outer average is not over
independently drawn spouses.

The same functional-independence argument applies to either component
of each vector. Their conditional means are $qP_M$ and $qP_F$, so
\begin{equation}
 C:=\mathbb E[X_MX_F^{\mathsf T}]
 =\mathbb E_{(P_M,P_F)}[(qP_M)(qP_F)^{\mathsf T}]
 =\rho qq^{\mathsf T}
 =\rho\begin{pmatrix}h^2&he\\he&e^2\end{pmatrix}.
 \label{eq:ge-cross}
\end{equation}
Thus $\operatorname{Cov}(E_M,E_F)=\rho e^2$ and
$\operatorname{Cov}(A_M,E_F)=\operatorname{Cov}(E_M,A_F)=\rho he$.
The full component covariance matrix is
$\left(\begin{smallmatrix}I_2&C\\C^{\mathsf T}&I_2\end{smallmatrix}\right)$.
Here the off-diagonal blocks are matrices, not scalars; $C$ happens to
be symmetric. Since $q^{\mathsf T}q=1$, $C$ has eigenvalue $\rho$ along $q$ and zero in the perpendicular direction. The full matrix therefore has eigenvalues
$1+\rho,1-\rho,1,1$, all positive in the stated parameter range.
Also $q^{\mathsf T}Cq=\rho(q^{\mathsf T}q)^2=\rho$, verifying the
specified spouse phenotype correlation. This is Otto et al.'s
equation~(22) after deleting the cultural component.

The intuition is that spouses match on a phenotype reflecting both
components. A high phenotype predicts a genetic value of $hP$, not
$P$ itself. Applying that prediction on both sides of the marriage
multiplies phenotype covariance by $h^2$. The relationship
$\rho_A=h^2\rho$ is therefore a \emph{derived restriction}, not a third
freely chosen parameter. Without (L), PH instead gives
$\operatorname{Cov}\{\mu(P_M),\mu(P_F)\}$, where
$\mu(p)=\mathbb E[A\mid P=p]$; it need not equal $h^2\rho$.

Although this restriction appears in the equilibrium model, its
derivation requires no intergenerational argument. Stationarity of
genetic variance is a separate condition, to which we now turn.

\subsection{Births and stationarity of genetic variance}
\label{sec:model-transmission}

For child $j$ of parents $M,F$, assume additive transmission:
\begin{equation}
 A_j=\tfrac12(A_M+A_F)+S_j,\qquad P_j=hA_j+eE_j.
 \label{eq:model-transmission}
\end{equation}
$S_j$ is the segregation deviation: the child's departure from the
midparental genetic value. The coefficient $1/2$ describes additive
Mendelian transmission; it does not specify the variance of $S_j$.

Let $\mathcal H$ collect the already realized ancestral and parental
components in the pedigree before these births. For distinct births,
assume
\[
 \mathbb E[S_j\mid\mathcal H]=0,\qquad
 \mathbb E[S_jS_k\mid\mathcal H]=0\quad(j\ne k).
\]
Draw offspring environmental values independently of one another,
of $\mathcal H$, and of the segregation deviations for these births,
with mean zero and variance one. These are sufficient
no-shared-environment and no-environmental-transmission assumptions.
They concern fresh offspring environments; they do not assert
independence from mates subsequently selected on offspring phenotypes.

For any mean-zero variable $Z$ with finite variance determined by $\mathcal H$,
\[
 \operatorname{Cov}(Z,S_j)
 =\mathbb E_{\mathcal H}\{Z\,\mathbb E[S_j\mid\mathcal H]\}=0,\qquad
 \operatorname{Cov}(S_j,S_k)
 =\mathbb E_{\mathcal H}\{\mathbb E[S_jS_k\mid\mathcal H]\}=0.
\]
Here the inner means use the conditional segregation distribution given
the existing family; the outer mean uses that family's joint
distribution. These equations explain why segregation cross-terms
vanish in what follows.

Write $W=(A_M+A_F)/2$. For an arbitrary spouse genetic correlation
$\rho_A$,
\begin{align*}
 \operatorname{Var}(A_j)
 &=\operatorname{Var}(W)+2\operatorname{Cov}(W,S_j)
   +\operatorname{Var}(S_j)\\
 &=\tfrac14\{\operatorname{Var}(A_M)+\operatorname{Var}(A_F)
              +2\operatorname{Cov}(A_M,A_F)\}
   +0+\operatorname{Var}(S_j)\\
 &=\tfrac12(1+\rho_A)+\operatorname{Var}(S_j).
\end{align*}
Stationarity, $\operatorname{Var}(A_j)=1$, therefore requires
\begin{equation}
 \boxed{\operatorname{Var}(S_j)=\tfrac12(1-\rho_A)
          =\tfrac12(1-h^2\rho).}
 \label{eq:model-segregation}
\end{equation}
Random mating is the special case $\rho=\rho_A=0$, giving variance
$1/2$; zero correlation by itself need not imply independent mating
outside a Gaussian specification. Fresh $E_j$ then implies
$\operatorname{Cov}(A_j,E_j)=0$ and $\operatorname{Var}(P_j)=1$.
This is a stationary moment specification, not a proof of convergence
of a multilocus population to equilibrium. Nor does it derive $h^2$
from allele frequencies or a random-mating genetic variance.

\subsection{Parent--child and full-sibling correlations}
\label{sec:model-nuclear}

All phenotype variances equal one, so phenotype covariances below
are correlations. Start with the mother and child. From the
within-person assumptions and the spouse block $C$,
\begin{align*}
 \operatorname{Cov}(P_M,A_M)&=h,\\
 \operatorname{Cov}(P_M,A_F)
 &=h\operatorname{Cov}(A_M,A_F)+e\operatorname{Cov}(E_M,A_F)\\
 &=h(\rho h^2)+e(\rho he)
   =\rho h(h^2+e^2)=\rho h.
\end{align*}
The second line is where the PH-derived spouse covariances enter.
Using transmission and the zero segregation covariance,
\[
 \operatorname{Cov}(P_M,A_j)
 =\tfrac12 h+\tfrac12\rho h+0=\tfrac12h(1+\rho).
\]
The child's fresh environment has zero covariance with $P_M$. Therefore
\begin{equation}
 \rho_{PO}:=\operatorname{Corr}(P_M,P_j)
 =h\operatorname{Cov}(P_M,A_j)+e\operatorname{Cov}(P_M,E_j)
 =\boxed{\tfrac12h^2(1+\rho).}
 \label{eq:po}
\end{equation}
Sex symmetry gives the same formula for the father.

To expose the environmental qualification explicitly,
\begin{align*}
 \operatorname{Cov}(A_M,A_j)
 &=\tfrac12\{1+\rho h^2\}+0,\\
 \operatorname{Cov}(E_M,A_j)
 &=\tfrac12\{\operatorname{Cov}(E_M,A_M)
             +\operatorname{Cov}(E_M,A_F)\}+0
   =\tfrac12\rho he.
\end{align*}
Consequently
\[
 \operatorname{Cov}(X_M,X_j)=
 \begin{pmatrix}(1+\rho h^2)/2&0\\\rho he/2&0\end{pmatrix}.
\]
The parent's environmental component is associated with the other
parent's genes through assortment, and those genes enter the child.
No parental environment is transmitted. Goldberger makes this
distinction explicitly in his Appendix, p.~A-5. It would be incorrect
to obtain parent--child resemblance by retaining only
$h^2\operatorname{Cov}(A_M,A_j)$.

For full siblings 1 and 2, $A_1=W+S_1$ and $A_2=W+S_2$, so
\begin{align*}
 \operatorname{Cov}(A_1,A_2)
 &=\operatorname{Var}(W)+\operatorname{Cov}(W,S_2)
   +\operatorname{Cov}(S_1,W)+\operatorname{Cov}(S_1,S_2)\\
 &=\operatorname{Var}(W)+0+0+0
   =\tfrac12(1+\rho h^2).
\end{align*}
Fresh environments give
$\operatorname{Cov}(A_1,E_2)=\operatorname{Cov}(E_1,A_2)
=\operatorname{Cov}(E_1,E_2)=0$. Expanding both phenotypes,
\begin{align}
 \rho_S:=\operatorname{Corr}(P_1,P_2)
 &=h^2\operatorname{Cov}(A_1,A_2)
   +he\operatorname{Cov}(A_1,E_2)\notag\\
 &\quad+eh\operatorname{Cov}(E_1,A_2)
   +e^2\operatorname{Cov}(E_1,E_2)\notag\\
 &=\boxed{\tfrac12h^2(1+\rho h^2).}
 \label{eq:sib}
\end{align}
These are the zero-cultural-transmission, zero-shared-environment
specializations of Otto et al.'s equations~(46) and~(55).
Adding sibling environmental covariance $c$ would add $e^2c$ if the
genetic--environmental cross-covariances remained zero. The classical
formula excludes that contribution by setting $c=0$.

\subsection{Extending the pedigree}
\label{sec:model-extended}

A single mating pair does not specify how a new spouse relates to the
rest of a pedigree. To derive extended-family moments we make that
extension explicit. For each existing person $i$, draw an outside mate
$\widetilde i$ so that the pair satisfies PH and the same marginal
assumptions. Given the existing person's phenotype, the mate's
components are independent of the rest of the already realized family.
Draw distinct outside mates independently conditional on that family.
Also assume the mate-phenotype conditional mean is linear:
\begin{equation}
 \mathbb E[P_{\widetilde i}\mid P_i]=\rho P_i.
 \tag{ML}\label{eq:model-mate-linearity}
\end{equation}
The coefficient is $\rho$ because both phenotype variances are one.
Neither PH nor (L) by itself implies (ML). These additional conditions
exclude other routes connecting an incoming spouse to the pedigree.
They supply a sufficient family sampling model for the following
formulas; they were not needed for the nuclear-family results.

Let $\mathcal F$ denote the existing family before these outside mates
and their children are drawn. PH, (L), and (ML) give
\begin{align*}
 \mathbb E[A_{\widetilde i}\mid\mathcal F]
 &=\mathbb E[A_{\widetilde i}\mid P_i]\\
 &=\mathbb E_{P_{\widetilde i}\mid P_i}
       \{\mathbb E[A_{\widetilde i}\mid P_{\widetilde i},P_i]\}\\
 &=\mathbb E_{P_{\widetilde i}\mid P_i}[hP_{\widetilde i}]
   =h\rho P_i.
\end{align*}
The first equality uses the family extension, the third uses the
PH conditioning reduction and (L), and the last uses (ML).
The displayed outer average is over the mate's phenotype conditional
on the existing person's phenotype. Hence for their child $c(i)$,
\begin{align}
 \mathbb E[A_{c(i)}\mid\mathcal F]
 &=\tfrac12\{A_i+h\rho P_i\}+0\notag\\
 &=\underbrace{\tfrac12(1+\rho h^2)}_{k}A_i
   +\underbrace{\tfrac12\rho he}_{d}E_i
   =kA_i+dE_i.
 \label{eq:model-child-mean}
\end{align}
The zero term follows by averaging the zero segregation mean first
over births and then over the incoming mate. For any centered
existing-family variable $Z$, the law of iterated expectations now gives
\[
 \operatorname{Cov}(Z,A_{c(i)})
 =\mathbb E_{\mathcal F}[Z(kA_i+dE_i)]
 =k\operatorname{Cov}(Z,A_i)+d\operatorname{Cov}(Z,E_i).
\]

For a grandparent $g$, their child $i$, and grandchild $c(i)$, the
parent--child calculation gives
$\operatorname{Cov}(P_g,A_i)=h(1+\rho)/2$, whereas the fresh environment
of $i$ gives $\operatorname{Cov}(P_g,E_i)=0$. Therefore
\begin{align*}
 \operatorname{Corr}(P_g,P_{c(i)})
 &=h\{k\operatorname{Cov}(P_g,A_i)
       +d\operatorname{Cov}(P_g,E_i)\}
       +e\operatorname{Cov}(P_g,E_{c(i)})\\
 &=h\{kh(1+\rho)/2+d\cdot0\}+0
   =\boxed{\tfrac12h^2(1+\rho)k}.
\end{align*}

For an aunt/uncle $u$ and their full sibling $i$,
$\operatorname{Cov}(A_u,A_i)=k$ and all cross-component environmental
covariances between the siblings are zero. Thus
$\operatorname{Cov}(P_u,A_i)=hk$ and
$\operatorname{Cov}(P_u,E_i)=0$. Their niece/nephew satisfies
\[
 \operatorname{Corr}(P_u,P_{c(i)})
 =h\{k(hk)+d\cdot0\}+e\cdot0=\boxed{h^2k^2}.
\]

For first cousins $c(i)$ and $c(u)$, independent outside-mate draws
and the birth assumptions make their genetic values conditionally
uncorrelated given $\mathcal F$. The law of total covariance
(the sum of the mean conditional covariance and the covariance of
the conditional means) therefore gives
\begin{align*}
 \operatorname{Cov}(A_{c(i)},A_{c(u)})
 &=\mathbb E_{\mathcal F}
       [\operatorname{Cov}(A_{c(i)},A_{c(u)}\mid\mathcal F)]\\
 &\quad+\operatorname{Cov}_{\mathcal F}(kA_i+dE_i,kA_u+dE_u)\\
 &=0+k^2\operatorname{Cov}(A_i,A_u)
   +kd\operatorname{Cov}(A_i,E_u)\\
 &\quad+dk\operatorname{Cov}(E_i,A_u)
   +d^2\operatorname{Cov}(E_i,E_u)\\
 &=k^2k+0+0+0=k^3.
\end{align*}
Both cousins have fresh environments, so their phenotype correlation is
$\boxed{h^2k^3}$. This calculation exposes exactly where the additional
assumption about different incoming mates is used.

\subsection{Model restrictions, estimation, and interpretation}
\label{sec:model-tests}

Table~\ref{tab:notation} collects the moments derived above. They
match the individual-relative rows of Clark's Technical Appendix,
Table~A1, with measurement error suppressed. In Clark's notation
$r=\rho$ and $m=\rho_A=h^2\rho$; his attenuation parameter is set to
one here. Goldberger's notation has $A=\rho_A$, $c_1=h^2$, and $c_2=1$
when dominance is omitted.
\begin{table}[htbp]
\centering
\caption{Predicted phenotype correlations; $k=(1+h^2\rho)/2$}
\label{tab:notation}
\begin{tabular}{ll}
\hline
Relationship & Correlation\\
\hline
Spouses & $\rho$\\
Parent--child & $h^2(1+\rho)/2$\\
Full siblings & $h^2k$\\
Grandparent--grandchild & $h^2(1+\rho)k/2$\\
Aunt/uncle--niece/nephew & $h^2k^2$\\
First cousins & $h^2k^3$\\
\hline
\end{tabular}
\end{table}

One directly testable implication is
\begin{equation}
 \rho_{PO}-\rho_S
 =\tfrac12h^2\{(1+\rho)-(1+h^2\rho)\}
 =\boxed{\tfrac12h^2\rho(1-h^2)>0}.
 \label{eq:gap}
\end{equation}
The strict inequality uses both positive assortment and
$0<h^2<1$. Under random mating, parent--child and full-sibling
correlations both reduce to $h^2/2$; grandparent and avuncular
correlations reduce to $h^2/4$, and first-cousin correlation to $h^2/8$.
Comparisons with estimated correlations must account for sampling
uncertainty and comparable measurement and sampling across kinships.

For estimation, the no-error model has two free structural parameters,
$\theta=(h^2,\rho)$, not three independent parameters
$(h^2,\rho,\rho_A)$. For population correlations $r_\ell$ of the
relationships in the table, write
\[
 r_\ell=f_\ell(h^2,\rho).
\]
A joint minimum-distance estimator, for example, minimizes
\[
 [\widehat{\boldsymbol r}-\boldsymbol f(h^2,\rho)]^{\mathsf T}
 W[\widehat{\boldsymbol r}-\boldsymbol f(h^2,\rho)]
\]
over admissible parameters, imposing $\rho_A=h^2\rho$ by substitution.
$\widehat{\boldsymbol r}$ is the vector of estimated kinship correlations;
$W$ is a positive-definite weighting matrix. A covariance-based choice
should account for dependence between correlation estimates from
overlapping families. This formulation constrains the entire collection
of moments jointly rather than choosing each relationship separately.

\citet[p.~19]{Goldberger1978} identifies precisely the error of treating
the spousal genetic correlation as free while retaining formulas
derived under Fisher's restriction. He also notes that alternative
mating specifications may remove that restriction but then require
rederiving the kinship formulas. A broader assortment trait therefore
cannot simply be invoked to free $\rho_A$ while leaving the formulary
unexamined.

A statistically established failure of $\rho_A=h^2\rho$ rejects the
joint assumptions underlying that restriction. It does not identify
which assumption failed. Small discrepancies among estimates can be
sampling error, and a formal comparison requires a coherent unrestricted
model and its sampling covariance. Under misspecification, fitted
parameters need not identify heritability or the claimed assortment
parameters; good fit supplies no guarantee against bias. Rejection
alone does not prove that every estimate is biased, or determine the
bias's direction or magnitude. Conversely, imposing the restriction is
necessary to estimate this particular model coherently, but does not
make its environmental exclusions empirically true.

Clark's presentation obscures this distinction. On p.~16 he questions
the implication that genetic resemblance must be weaker than phenotype
resemblance for educational attainment, appealing to matching on a wider
set of characteristics. On p.~75 he instead invokes an imperfectly
measured underlying phenotype and notes that the three parameters are
linked. The latter acknowledgment does not resolve the former departure
from the model. The issue is not whether the relationship appears
somewhere in the book, but whether it is consistently treated as a
restriction required by the assumptions generating the fitted equations.

Two different arguments need to be separated. If years of education
is merely a noisy measure of the phenotype on which spouses assort,
then $\rho_A=h^2\rho$ still applies to that underlying phenotype:
$\rho$ must refer to its spousal correlation, and a measurement model
must connect its moments to those observed. If spouses instead match
on several characteristics in a way that violates PH for the modeled
phenotype, the matching model has changed. That possibility is legitimate,
but it requires specifying the new joint distribution and deriving the
resulting kinship moments. The book's appeal to broader matching does
not supply that specification or demonstrate that the same kinship
formulas remain valid. The restriction cannot be treated as an optional
extra while retaining the interpretation of formulas derived using it.

This leaves a gap between the fit and the substantive conclusion.
Clark interprets the pattern as evidence for additive genetic
transmission combined with very strong assortment on underlying social
abilities \citep[pp.~16,75,286 and Technical Appendix]{Clark2026}.
But a close fit of the correlation formulas does not, by itself,
demonstrate that this genetic mechanism---or assortment on a particular
social or cognitive trait---generated the data. That interpretation
requires a coherent specification connecting the matching process,
the transmitted components, and the observed outcomes, with its
restrictions imposed and tested. Moving between the classical model
and informal alternatives leaves the reader without a clear account
of which model the reported fit actually supports.

\paragraph{Verification.}
The accompanying verification script checks the spouse and family
covariance matrices algebraically and simulates an explicit Gaussian
construction satisfying the assumptions. It checks parent--child,
full-sibling, grandparent, avuncular, and first-cousin moments, as well
as component covariances and stationary variances. All 46 symbolic checks passed. Across nine combinations of $h^2$ and
$\rho$, simulations of 2.25 million independent pedigrees checked 1,224
component covariances and 54 phenotype moments. The largest absolute
standardized Monte Carlo discrepancies were 3.402 and 2.091, respectively,
within the prespecified six-standard-error diagnostic threshold.
The script is \texttt{verify\_fisher\_model.py}, in the repository's
\texttt{verification/} directory. The construction
provides an existence and calculation check for this moment model;
simulation is not a proof of empirical adequacy or of multilocus
equilibrium.
